\documentclass[aspectratio=169]{beamer}

\IfFileExists{beamerthememetropolis.sty}{\usetheme{metropolis}}{\usetheme{default}}
\usepackage[utf8]{inputenc}
\usepackage[T1]{fontenc}
\usepackage{booktabs}
\usepackage{graphicx}
\usepackage{amsmath}
\usepackage{xcolor}

\definecolor{shipblue}{RGB}{20,74,130}
\definecolor{warningorange}{RGB}{210,105,30}
\setbeamercolor{structure}{fg=shipblue}
\setbeamercolor{alerted text}{fg=warningorange}
\setbeamertemplate{navigation symbols}{}

\newcommand{\results}{../../results/scan_hi_2026-06-08}

\title{SHiP Timing Detector\\Estado de la simulación Geant4}
\author{René Ríos (ULS)}
\institute{Revisión de supervisor / SHiP Collaboration Meeting}
\date{Junio 2026 \quad Geant4 11.4.0}

\begin{document}

\begin{frame}
  \titlepage
\end{frame}

\begin{frame}{Detector y scan de alta estadística}
  \begin{columns}[T]
    \column{0.55\textwidth}
    \begin{itemize}
      \item Barra activa EJ-204: $1400\times60\times10$ mm$^3$.
      \item End-left: IDs 0--7; End-right: IDs 8--15.
      \item Top: IDs 16--35 sobre la cara superior.
      \item Comparación: \textbf{Side/End readout} frente a \textbf{Top readout}.
    \end{itemize}
    \column{0.42\textwidth}
    \begin{block}{Configuración reproducible}
      $\mu^-$ de 1 GeV, incidencia vertical. 31 posiciones:
      paso de 50 mm entre $-650$ y $+650$ mm, más $\pm670$ y $\pm690$ mm.
      10\,000 eventos y semillas fijas por posición.
    \end{block}
  \end{columns}
\end{frame}

\begin{frame}{Óptica y punto abierto del wrapping}
  \begin{itemize}
    \item Modelo actual: panel reflector explícito, pulido, $R(\lambda)=0.98$.
    \item Las superficies de los SiPM permanecen separadas del reflector.
    \item Montaje físico por cerrar: Tyvek (principalmente difuso/Lambertiano)
      frente a Mylar aluminizado (más especular).
  \end{itemize}
  \begin{alertblock}{Impacto pendiente}
    La elección difuso/especular altera la dispersión de caminos, el yield y
    la resolución temporal; no se modificó en este scan.
  \end{alertblock}
\end{frame}

\begin{frame}{Baseline física aplicada}
  \centering
  \begin{tabular}{lrrrrr}
    \toprule
    Material & Yield & $\tau_r$ & $\tau_d$ & Aten. & $\lambda_{\max}$\\
             & ph/MeV & ns & ns & cm & nm\\
    \midrule
    EJ-204 & 10400 & 0.7 & 1.8 & 160 & 408\\
    EJ-200 & 10000 & 0.9 & 2.1 & 380 & 425\\
    EJ-230 & 9700  & 0.5 & 1.5 & 120 & 391\\
    \bottomrule
  \end{tabular}
  \vspace{0.5cm}
  \begin{itemize}
    \item Catálogo por datasheet; corrección EJ-204: atenuación = \textbf{160 cm}.
    \item EJ-204 activo por defecto, seleccionable y registrado al inicio del run.
    \item Guardrail CTest \texttt{physics\_baseline\_check}.
  \end{itemize}
\end{frame}

\begin{frame}{Rise time finito y CFD@14\%}
  \[
    I(t) \propto e^{-t/\tau_d}-e^{-t/\tau_r}, \qquad t>0
  \]
  \begin{itemize}
    \item Antes: emisión instantánea en $t_0$, con flanco temprano irreal.
    \item Ahora: perfil biexponencial con $(\tau_r,\tau_d)$ del datasheet.
    \item El flanco temprano fija el cruce CFD@14\%; por ello el rise time
      efectivo es crítico para una $\sigma_t$ defendible.
    \item Sin waveform, CFD@14\% se aproxima por el fotoelectrón
      $k=\mathrm{round}(0.14N_{\rm pe})$; FPT queda como cross-check.
    \item End usa $(t_L+t_R)/2$ y exige al menos 10 PE en cada lado.
  \end{itemize}
\end{frame}

\begin{frame}{Yield y validación geométrica}
  \centering
  \includegraphics[width=0.88\textwidth]{\results/yield_vs_position.pdf}

  \scriptsize El intercambio End-L/End-R bajo $x\to-x$ valida la simetría
  longitudinal; Top permanece aproximadamente independiente de la posición.
\end{frame}

\begin{frame}{Resolución intrínseca CFD@14\% plegada}
  \centering
  \includegraphics[width=0.75\textwidth]{\results/sigma_cfd_folded_vs_abs_position.pdf}
  \begin{alertblock}{Caveat central}
    \small $\sigma_t$ es \textbf{intrínseco}: SPTR y jitter electrónico son cero.
    Es una cota inferior, todavía no comparable con 250 ps del test beam.
    Los puntos End con un lado bajo 10 PE se marcan como no fiables.
  \end{alertblock}
\end{frame}

\begin{frame}{Validación de simetría tras alta estadística}
  \begin{columns}[T,onlytextwidth]
    \column{0.47\textwidth}
    \centering
    \includegraphics[width=\textwidth]{\results/sigma_cfd_unfolded_vs_position.pdf}
    \column{0.50\textwidth}
    \scriptsize
    \begin{block}{Métrica}
      \[
        A_\sigma(|x|)=
        \frac{2[\sigma(+x)-\sigma(-x)]}
             {\sigma(+x)+\sigma(-x)}
      \]
      Se evalúa sin plegar; las curvas finales de $\sigma$ combinan $\pm|x|$.
    \end{block}
    \begin{itemize}
      \item Yield End-L/End-R simétrico: geometría consistente.
      \item La asimetría residual de $\sigma$ cae frente al scan de 200 eventos.
      \item Cualquier residuo significativo persistente apunta al análisis,
        no a la física simétrica.
    \end{itemize}
    \begin{block}{Resultado}
      \small End fiable: mediana $|A_\sigma|=2.35\%$, máximo $4.15\%$.
      Top: mediana $|A_\sigma|=1.12\%$. La asimetría del scan preliminar
      estaba dominada por estadística.
    \end{block}
  \end{columns}
\end{frame}

\begin{frame}{Resumen CFD@14\% en la región End fiable}
  \centering
  \begin{tabular}{rrrrr}
    \toprule
    $|x|$ [mm] & $N_{\rm eff,End}$ & $\sigma_{\rm End}$ [ps] &
    $\sigma_{\rm Top}$ [ps] & End/Top\\
    \midrule
      0 &  9662 & $186.4\pm2.2$ & $76.3\pm0.9$ & 2.44\\
     50 & 18724 & $188.3\pm1.6$ & $73.3\pm0.6$ & 2.57\\
    100 & 16950 & $187.4\pm1.7$ & $70.9\pm0.6$ & 2.64\\
    150 & 14441 & $186.5\pm1.9$ & $76.2\pm0.6$ & 2.45\\
    200 & 11224 & $192.0\pm2.2$ & $75.5\pm0.6$ & 2.54\\
    \bottomrule
  \end{tabular}
  \vspace{0.5cm}
  \begin{block}{Resultado}
    En la región fiable, Top mejora $\sigma_t$ intrínseco por un factor
    mediano de \textbf{2.54}; rango 2.44--2.64. Los errores son estadísticos
    del ajuste gaussiano iterativo del core.
  \end{block}
\end{frame}

\begin{frame}{Cross-check FPT frente a CFD@14\%}
  \centering
  \includegraphics[width=0.86\textwidth]{\results/sigma_fpt_vs_cfd_folded.pdf}

  \scriptsize CFD@14\% es el estimador primario. FPT conserva continuidad con
  el scan preliminar, pero es más sensible a los primeros arribos y presenta
  colas no gaussianas.
\end{frame}

\begin{frame}{Eficiencia de detección}
  \centering
  \includegraphics[width=0.84\textwidth]{\results/efficiency_vs_position.pdf}

  \scriptsize Definición del proyecto:
  $100\,N_{\mathrm{detectados}}/N_{\mathrm{centelleo\ generados}}$.
\end{frame}

\begin{frame}{Lectura física}
  \begin{itemize}
    \item \textbf{Top readout}: SiPMs a milímetros de la traza capturan fotones
      antes de acumular dispersión por propagación; yield relativamente estable
      y $\sigma_{\rm Top}<\sigma_{\rm End}$ a nivel intrínseco.
    \item \textbf{End readout}: atenuación y dispersión de caminos degradan el
      lado lejano; cuando queda hambriento, $(t_L+t_R)/2$ deja de ser fiable.
    \item \textbf{Trade-off EJ-204}: atenuación de 160 cm frente a 380 cm de
      EJ-200 produce yields End más bajos y empinados, a cambio de centelleo
      más rápido ($0.7/1.8$ ns).
  \end{itemize}
  \begin{block}{Titular defendible}
    Top mejora de forma robusta la resolución temporal intrínseca frente a End
    por un factor mediano de 2.54 en la región fiable. Los valores absolutos
    quedan pendientes del jitter real.
  \end{block}
\end{frame}

\begin{frame}{Próximos pasos y puntos abiertos}
  \begin{itemize}
    \item EXEC\_02b: SPTR $\oplus$ jitter electrónico en el ToA, tras confirmar
      el modelo Hamamatsu y la procedencia del SPTR.
    \item Emular FastIC+: 1 LSB = 250 ps, Type 1/2, cortes $\geq300/\geq500$ LSB.
    \item Decidir y calibrar reflector difuso frente a especular.
    \item Aclarar la inconsistencia del factor $1/\sqrt{8}$.
  \end{itemize}
  \begin{block}{Composición futura}
    \[
      \sigma_{\rm total} =
      \sqrt{\sigma_{\rm intr}^2+\sigma_{\rm SPTR}^2+\sigma_{\rm elec}^2}
    \]
  \end{block}
\end{frame}

\end{document}
